\chapter{Cadre général du projet}

% ============================================================
\section{Introduction}
% ============================================================

Ce premier chapitre pose le cadre dans lequel s'inscrit notre projet de fin d'études.
Nous présentons successivement l'organisme d'accueil TAC-TIC, la problématique
identifiée, l'étude critique des solutions existantes, la solution proposée et
la méthodologie de développement retenue.

% ============================================================
\section{Présentation de l'organisme d'accueil}
% ============================================================

\subsection{Présentation générale}

\textbf{TAC-TIC}, filiale du Groupe New Ways, est une entreprise de services
numériques (ESN) fondée en 2012 et implantée au Technoparc El Ghazela, Ariana
(Tunisie). Elle compte plus d'une centaine de collaborateurs et est certifiée
\textbf{ISO 27001} (sécurité de l'information) et \textbf{ISO 9001} (management
de la qualité). Son cœur de métier couvre le développement web et mobile,
l'architecture cloud et l'intégration de solutions d'intelligence artificielle
et d'IoT.

\begin{center}
  \IfFileExists{figures/logo.png}{%
    \includegraphics[width=0.30\textwidth]{figures/logo.png}%
  }{%
    \fbox{\parbox{0.30\textwidth}{\centering\vspace{1cm}
      \textbf{[Logo TAC-TIC]}\\\texttt{figures/logo.png}
    \vspace{1cm}}}%
  }
\end{center}

\begin{table}[H]
  \centering
  \caption{Jalons historiques de TAC-TIC}
  \label{tab:historique_tactic}
  \begin{tabularx}{\textwidth}{C{2cm}X}
    \toprule
    \thead{Année} & \thead{Étape clé} \\
    \midrule
    \rowcolor{lightgray}
    2012 & Création de TAC-TIC au sein du Groupe New Ways ; premiers projets web \\
    2014 & Obtention de la certification ISO 27001 ; spécialisation en cybersécurité \\
    \rowcolor{lightgray}
    2017 & Transition vers l'architecture cloud (AWS, Azure) \\
    2020 & Lancement de la division IA/IoT ; déploiement de premiers modèles ML en production \\
    \rowcolor{lightgray}
    2023 & Dépassement du cap des 100 projets livrés ; ouverture d'un bureau à Paris \\
    2026 & Développement de TAC-TIC IntelliPM (présent PFE) \\
    \bottomrule
  \end{tabularx}
\end{table}

\begin{table}[H]
  \centering
  \caption{Domaines d'activité de TAC-TIC}
  \label{tab:domaines_tactic}
  \begin{tabularx}{\textwidth}{L{3.5cm}X}
    \toprule
    \thead{Domaine} & \thead{Activités principales} \\
    \midrule
    \rowcolor{lightgray}
    Développement web     & Applications SPA React/Vue, APIs REST Laravel/Node.js \\
    Développement mobile  & Applications iOS/Android (React Native, Flutter) \\
    \rowcolor{lightgray}
    Cloud \& DevOps       & Architecture cloud, CI/CD, Docker, monitoring \\
    Intelligence artificielle & ML supervisé/non supervisé, NLP, vision par ordinateur \\
    \rowcolor{lightgray}
    Conseil \& audit      & Audit sécurité ISO 27001, conseil en transformation numérique \\
    \bottomrule
  \end{tabularx}
\end{table}

\subsection{Service d'affectation}

Notre stage s'est déroulé au sein du \textbf{département Développement informatique},
division web/cloud, encadrés par le responsable technique de la division.
Ce département réalise les projets de développement logiciel pour les clients
internes et externes de TAC-TIC, en adoptant des méthodologies agiles et des
architectures distribuées.

\subsection{Organigramme}

\begin{figure}[H]
  \centering
  \IfFileExists{figures/organigramme_tactic.png}{%
    \includegraphics[width=0.85\textwidth]{figures/organigramme_tactic.png}%
  }{%
    \fbox{\parbox{0.85\textwidth}{\centering\vspace{1.5cm}
      \textbf{[IMAGE À INSÉRER]}\\[0.4em]
      \texttt{figures/organigramme\_tactic.png}\\[0.3em]
      \textit{Organigramme de l'entreprise TAC-TIC — Groupe New Ways}
    \vspace{1.5cm}}}%
  }
  \caption{Organigramme de TAC-TIC — Groupe New Ways}
  \label{fig:organigramme}
\end{figure}

% ============================================================
\section{Présentation du projet}
% ============================================================

\subsection{Contexte}

TAC-TIC gère simultanément une dizaine à une quinzaine de projets actifs,
mobilisant des équipes de 3 à 12 développeurs par projet. La coordination de ces
équipes, la traçabilité des décisions et le respect des délais constituent des défis
quotidiens. Les outils utilisés jusqu'alors — combinaison de Trello et de tableurs
Excel — ne permettaient pas de disposer d'une vue consolidée multi-projets,
d'une planification fine par sprints, ni d'une aide à la décision par intelligence
artificielle pour l'assignation et la gestion des risques.

\subsection{Problématique}

\begin{mdframed}[linecolor=primary,linewidth=1.5pt]
\centering
\textit{Comment concevoir une plateforme de gestion de projets qui intègre
nativement l'intelligence artificielle pour assister les chefs de projet dans
leurs décisions d'assignation, de planification et de gestion des risques,
tout en garantissant une architecture logicielle maintenable et une expérience
utilisateur fluide ?}
\end{mdframed}

\subsection{Objectifs généraux}

\begin{itemize}
  \item Centraliser la gestion de tous les projets sur une plateforme unique ;
  \item Adopter les cérémonies Scrum (sprints, rétrospective, vélocité) ;
  \item Offrir des visualisations décisionnelles (Gantt, Kanban, dashboards) ;
  \item Automatiser les alertes de retard et les recommandations d'assignation.
\end{itemize}

\subsection{Objectifs spécifiques}

\begin{itemize}
  \item Implémenter une architecture distribuée (backend Laravel + SPA React
        + microservice ML FastAPI) ;
  \item Mettre en place une sécurité multicouche (Sanctum, RBAC à 3 niveaux,
        policies, journal d'audit) ;
  \item Délivrer des événements temps réel via WebSocket (Laravel Reverb) ;
  \item Atteindre le Quality Gate SonarCloud (0 bug critique, duplication
        $\leq$ 3\,\%).
\end{itemize}

% ============================================================
\section{Étude de l'existant}
% ============================================================

\subsection{Solutions analysées}

Nous avons étudié trois outils majeurs du marché : \textbf{Trello}, \textbf{Jira}
et \textbf{Asana}.

\subsection{Comparaison des solutions existantes}

\begin{table}[H]
  \centering
  \caption{Comparaison des solutions de gestion de projets existantes}
  \label{tab:comparaison_existant}
  \small
  \begin{tabularx}{\textwidth}{L{4.5cm}C{2cm}C{2cm}C{2cm}C{2.5cm}}
    \toprule
    \thead{Fonctionnalité} & \thead{Trello} & \thead{Jira} & \thead{Asana}
                           & \thead{TAC-TIC IntelliPM} \\
    \midrule
    \rowcolor{lightgray}
    Kanban drag-and-drop          & \cmark & \cmark & \cmark & \cmark \\
    Diagramme de Gantt            & \xmark & \cmark & \cmark & \cmark \\
    \rowcolor{lightgray}
    Chemin critique (CPM)         & \xmark & \xmark & \xmark & \cmark \\
    Gestion des sprints           & \xmark & \cmark & \cmark & \cmark \\
    \rowcolor{lightgray}
    Recommandation ML assignation & \xmark & \xmark & \xmark & \cmark \\
    Assistant IA conversationnel  & \xmark & \xmark & \xmark & \cmark \\
    \rowcolor{lightgray}
    Notifications WebSocket       & \xmark & \cmark & \cmark & \cmark \\
    RBAC granulaire (atomique)    & \xmark & partiel & \xmark & \cmark \\
    \rowcolor{lightgray}
    Journal d'audit enrichi       & \xmark & \xmark & \xmark & \cmark \\
    Open source / hébergeable     & \xmark & \xmark & \xmark & \cmark \\
    \bottomrule
  \end{tabularx}
\end{table}

\subsection{Critique des solutions existantes}

\begin{itemize}
  \item \textbf{Trello} : simple et visuel, mais limité aux tableaux Kanban.
        Aucune notion de sprint, pas de Gantt, pas d'IA.
  \item \textbf{Jira} : puissant mais complexe à configurer. Coûteux pour les
        PME. Aucune recommandation ML native.
  \item \textbf{Asana} : bon équilibre fonctionnel, mais pas de chemin critique,
        pas d'IA intégrée, coût élevé par siège.
\end{itemize}

\subsection{Positionnement d'IntelliPM}

TAC-TIC IntelliPM se différencie par l'intégration native d'un moteur
d'intelligence artificielle hybride (Machine Learning pour la recommandation
d'assignation + LLM pour l'assistance conversationnelle), d'un journal d'audit
enterprise et d'un système RBAC à permissions atomiques.

\begin{table}[H]
  \centering
  \caption{Apports différenciants de TAC-TIC IntelliPM}
  \label{tab:apports}
  \begin{tabularx}{\textwidth}{L{4.5cm}X}
    \toprule
    \thead{Fonctionnalité} & \thead{Apport} \\
    \midrule
    \rowcolor{lightgray}
    Recommandation ML       & Suggère le meilleur assigné via un modèle
                              GBM/Random Forest entraîné sur les données projet \\
    Assistant LLM           & Chat conversationnel avec contexte projet chargé
                              dynamiquement (OpenRouter / Llama 3.3 70B) \\
    \rowcolor{lightgray}
    Estimation de durée     & Calcul heuristique (story points, priorité, tâches
                              similaires) enrichi par LLM \\
    Détection de risques    & Score de risque par règles métier (délai, dépendances,
                              priorité) \\
    \rowcolor{lightgray}
    Chemin critique (CPM)   & Calcul automatique et visualisation sur le Gantt \\
    Temps réel WebSocket    & 11 événements broadcastés via Laravel Reverb \\
    \rowcolor{lightgray}
    RBAC enterprise         & Rôles + 35 permissions atomiques + rôles projet \\
    Journal d'audit         & Valeurs avant/après, IP, sévérité automatique,
                              export, broadcast WebSocket \\
    \bottomrule
  \end{tabularx}
\end{table}

% ============================================================
\section{Méthodologie adoptée}
% ============================================================

\subsection{Scrum}

La méthode \textbf{Scrum} a été retenue pour trois raisons principales :
\begin{enumerate}
  \item \textbf{Livraisons incrémentales} — chaque sprint livre un incrément
        fonctionnel démontrable, facilitant les retours du tuteur entreprise ;
  \item \textbf{Adaptabilité} — les backlogs sont révisés entre sprints pour
        intégrer les retours utilisateurs ;
  \item \textbf{Transparence} — les daily stand-ups et rétrospectives assurent
        une visibilité permanente sur l'avancement.
\end{enumerate}

\begin{figure}[H]
  \centering
  \IfFileExists{figures/cycle_scrum.png}{%
    \includegraphics[width=0.80\textwidth]{figures/cycle_scrum.png}%
  }{%
    \fbox{\parbox{0.80\textwidth}{\centering\vspace{1.5cm}
      \textbf{[IMAGE À INSÉRER]}\\[0.4em]
      \texttt{figures/cycle\_scrum.png}\\[0.3em]
      \textit{Cycle Scrum : Product Backlog → Sprint Backlog → Sprint → Incrément}
    \vspace{1.5cm}}}%
  }
  \caption{Cycle Scrum appliqué au projet TAC-TIC IntelliPM}
  \label{fig:cycle_scrum}
\end{figure}

\subsection{Organisation des sprints}

Le projet a été réalisé sur \textbf{douze semaines} réparties en six sprints
de deux semaines chacun.

\begin{table}[H]
  \centering
  \caption{Planification des 6 sprints du projet}
  \label{tab:sprints_planning}
  \begin{tabularx}{\textwidth}{C{1.5cm}L{5cm}X C{2.5cm}}
    \toprule
    \thead{Sprint} & \thead{Périmètre} & \thead{Livrables principaux} & \thead{Durée} \\
    \midrule
    \rowcolor{lightgray}
    1 & Authentification \& utilisateurs  & Auth Sanctum, RBAC Spatie, profil, dashboard admin & 2 semaines \\
    2 & Projets et équipes               & CRUD projets, invitations, sprints, équipes         & 2 semaines \\
    \rowcolor{lightgray}
    3 & Tâches, Kanban et Gantt          & Tâches hiérarchiques, dépendances, Kanban, Gantt CPM & 2 semaines \\
    4 & Notifications et collaboration   & WebSocket, calendrier, dashboards, recherche         & 2 semaines \\
    \rowcolor{lightgray}
    5 & Intelligence artificielle        & ML recommandation, estimation, risques, chat LLM     & 2 semaines \\
    6 & Administration et audit          & RBAC atomique, journal d'audit, analytics, PDF       & 2 semaines \\
    \bottomrule
  \end{tabularx}
\end{table}

% ============================================================
\section{Conclusion}
% ============================================================

Ce premier chapitre a délimité le cadre complet du projet : une entreprise
dynamique (TAC-TIC / Groupe New Ways), une problématique réelle de coordination
multi-projets, et une réponse technologique ambitieuse articulant agilité Scrum,
architecture distribuée et intelligence artificielle. Le chapitre suivant entre
dans le détail de l'analyse, de la planification et de la conception architecturale.
